% Part: normal-modal-logic
% Chapter: filtrations
% Section: examples-of-filtrations

\documentclass[../../../include/open-logic-section]{subfiles}

\begin{document}

\olfileid{nml}{fil}{exf}

\olsection{过滤的例子}

我们尚未证明是否存在过滤。但实际上，对于任意模型 $\mModel{M}$，都存在许多 $\mModel{M}$ 通过 $\Gamma$ 的过滤。我们特别指出两种：最细过滤和最粗过滤。同一模型的不同过滤，其区别在于可达关系（因为 \olref[fil]{defn:filtration} 直接规定了 $W^*$ 和~$V^*$ 应如何设定）。最细过滤将拥有尽可能少的可达世界，而最粗过滤则拥有尽可能多的可达世界。

\begin{defn}
  设 $\Gamma$ 对子公式封闭，模型 $\mModel{M}$ 的\emph{最细}过滤 $\mModel{M^*}$ 定义为规定：\[
  R^*[u][v] \quad \text{当且仅当} \quad \exists u'\in [u] \;
  \exists v' \in [v] : Ru'v'.
  \]
\end{defn}

\begin{prop}\ollabel{prop:finest}
  最细过滤 $\mModel{M^*}$ 确实是一个过滤。
\end{prop}

\begin{proof}
  我们需要检查如此定义的 $R^*$ 满足
  \olref[fil]{defn:filtration}\olref[fil]{defn:filtration-R}。我们依次检查这三个条件。

  如果 $Ruv$，那么由于 $u \in [u]$ 且 $v \in [v]$，也有 $R^*[u][v]$，
  因此满足 \olref[fil]{defn:filtration-R1}。

  \iftag{prvBox}{\iftag{probBox}{我们将
      \olref[fil]{defn:filtration-R2} 的验证留作练习。}{对于
        \olref[fil]{defn:filtration-R2}，假设 $\Box!A \in \Gamma$，
        $R^*[u][v]$，且 $\mSat{M}{\Box!A}[u]$。根据 $R^*$ 的定义，
        存在 $u' \equiv u$ 和 $v' \equiv v$ 使得 $Ru'v'$。
        由于 $u$ 和 $u'$ 在 $\Gamma$ 上一致，故也有
        $\mSat{M}{\Box!A}[u']$，从而 $\mSat{M}{!A}[v']$。
        由于 $\Gamma$ 对子公式封闭，$v$ 和 $v'$ 在 $!A$ 上一致，
        所以 $\mSat{M}{!A}[v]$，如所欲证。}}{}

  \iftag{prvDiamond}{\iftag{probDiamond}{我们将
      \olref[fil]{defn:filtration-R3} 的验证留作练习。}{为验证
        \olref[fil]{defn:filtration-R3}，假设 $\Diamond!A \in
        \Gamma$，$R^*[u][v]$，且 $\mSat{M}{!A}[v]$。
        根据 $R^*$ 的定义，存在 $u' \equiv u$ 和 $v' \equiv v$ 使得
        $Ru'v'$。由于 $v$ 和 $v'$ 在~$\Gamma$ 上一致，且 $\Gamma$
        对子公式封闭，故也有 $\mSat{M}{!A}[v']$，从而
        $\mSat{M}{\Diamond !A}[u']$。因为 $u$ 和 $u'$ 也在 $\Gamma$ 上一致，
        故得 $\mSat{M}{\Diamond !A}[u]$。}}{}
\end{proof}

\begin{probtag}{probBox,probDiamond}
  完成 \olref[nml][fil][exf]{prop:finest} 的证明。
\end{probtag}

\begin{defn}
  设 $\Gamma$ 对子公式封闭，模型~$\mModel{M}$ 的\emph{最粗}过滤~$\mModel{M^*}$
  通过以下方式定义：规定 $R^*[u][v]$ 当且仅当
  \iftag{notprvBox,notprvDiamond}{满足以下条件：}{\emph{同时}满足以下两个条件：}
  \begin{tagenumerate}{prvBox,prvDiamond}
  \tagitem{prvBox}{\ollabel{defn:coarsest-Box}若 $\Box!A \in \Gamma$
    且 $\mSat{M}{\Box!A}[u]$，则
    $\mSat{M}{!A}[v]$\iftag{prvDiamond}{；}{。}}{}
  \tagitem{prvDiamond}{\ollabel{defn:coarsest-Diamond}若 $\Diamond!A
    \in \Gamma$ 且 $\mSat{M}{!A}[v]$，则
    $\mSat{M}{\Diamond!A}[u]$。}{}
  \end{tagenumerate}
\end{defn}

\begin{prop}
  最粗过滤 $\mModel{M^*}$ 确实是一个过滤。
\end{prop}

\begin{proof}
  根据 $R^*$ 的定义，唯一尚需验证的条件是由 $Ruv$ 可推出 $R^*[u][v]$。
  故假设 $Ruv$。
  \iftag{prvBox}{假设 $\Box!A \in \Gamma$ 且
    $\mSat{M}{\Box!A}[u]$；那么显然有
    $\mSat{M}{!A}[v]$\iftag{prvDiamond}{，且
      满足 \olref{defn:coarsest-Box}。}{。}}{}
  \iftag{prvDiamond}{假设 $\Diamond!A \in
    \Gamma$ 且 $\mSat{M}{!A}[v]$。则由~$Ruv$ 可得 $\mSat{M}{\Diamond!A}[u]$
    \iftag{prvBox}{，且满足 \olref{defn:coarsest-Diamond}}{}。}{}
\end{proof}

\begin{ex}
  令 $W = \PosInt$，$Rnm$ 当且仅当 $m = n + 1$，并且 $V(p) = \Setabs{2n}{n \in
    \Nat}$。模型 $\mModel{M} = \tuple{W, R, V}$ 如
  \olref{fig:ex-filtration} 所示。这些世界是 $1$、$2$ 等等；每个世界恰好能到达另一个世界——它的后继——并且 $p$ 在且仅在偶数处为真。

  \begin{figure}
    \centering
    \begin{tikzpicture}[modal]
      \node[world] (1) [label=below:\mFalse{p}] {$1$};
      \node[world] (2) [label=below:\mTrue{p}, right=of 1] {$2$};
      \node[world] (3) [label=below:\mFalse{p}, right=of 2] {$3$};
      \node[world] (4) [label=below:\mTrue{p}, right=of 3] {$4$};
      \node[phantom] (5) [right=of 4] {};
      \draw[->] (1) to (2);
      \draw[->] (2) to (3);
      \draw[->] (3) to (4);
      \draw[dotted] (4) to (5);
    \end{tikzpicture}

    \begin{tikzpicture}[modal]
      \node[world] (1) [label=below:\mFalse{p}] {$[1]$};
      \node[world] (2) [label=below:\mTrue{p}, right=of 1] {$[2]$};
      \draw[->,bend left] (1) to (2);
      \draw[->,bend left] (2) to (1);

      \node[world] (3) [label=below:\mFalse{p},right=of 2] {$[1]$};
      \node[world] (4) [label=below:\mTrue{p}, right=of 3] {$[2]$};
      \draw[->,bend left] (3) to (4);
      \draw[->,bend left] (4) to (3);
      \draw[->,reflexive above] (4) to (4);
    \end{tikzpicture}
    \caption{一个无限模型及其过滤。}
    \ollabel{fig:ex-filtration}
  \end{figure}

  现在令 $\Gamma$ 为~$\Box p \lif p$ 的子公式集合，
  即 $\{p, \Box p, \Box p \lif p\}$。$p$ 在且仅在偶数处为真，
  $\Box p$ 在且仅在奇数处为真，
  所以 $\Box p \lif p$ 在且仅在偶数处为真。
  换言之，每个奇数使得 $\Box p$ 为真，但使得 $p$ 和 $\Box p \lif p$ 为假；
  每个偶数使得 $p$ 和 $\Box p \lif p$ 为真，但使得 $\Box p$ 为假。
  因此 $W^* = \{ [1], [2] \}$，其中 $[1] =
  \{1, 3, 5, \dots\}$ 且 $[2] = \{2, 4, 6, \dots\}$。
  由于 $2 \in V(p)$，有 $[2] \in V^*(p)$；
  由于 $1 \notin V(p)$，有 $[1] \notin V^*(p)$。
  所以 $V^*(p) = \{[2]\}$。

  任何基于 $W^*$ 的过滤，其可达关系必须包含 $\tuple{[1], [2]}, \tuple{[2],[1]}$：
  由于 $R12$，根据
  \olref[fil]{defn:filtration}\olref[fil]{defn:filtration-R1} 我们必须有 $R^*[1][2]$；
  并且由于 $R23$ 我们必须有 $R^*[2][3]$，而 $[3]=[1]$。
  它不能包含 $\tuple{[1],[1]}$：若是包含，我们将有 $R^*[1][1]$，
  $\mSat{M}{\Box p}[1]$ 但 $\mSat/{M}{p}[1]$，这与
  \olref[fil]{defn:filtration-R2} 矛盾。
  没有任何条件要求或排除 $R^*[2][2]$。
  因此，~$\mModel{M}$ 有两种可能的过滤，分别对应两种可达关系
  \[
  \{\tuple{[1],[2]}, \tuple{[2],[1]}\} \text{ 和 }
  \{\tuple{[1],[2]}, \tuple{[2],[1]}, \tuple{[2],[2]}\}.
  \]
  在任何一种情况下，在~$[1]$ 处 $p$ 和 $\Box p \lif p$ 为假而 $\Box p$ 为真；
  在~$[2]$ 处 $p$ 和 $\Box p \lif p$ 为真而 $\Box p$ 为假。
\end{ex}

\begin{prob}
  考虑如下模型 $\mModel{M} = \tuple{W, R, V}$，其中 $W
  = \Setabs{0\sigma}{\sigma \in \Bin^*}$ 是由 $0$ 和~$1$ 组成且以~$0$ 开头的序列的集合，$R\sigma\sigma'$ 当且仅当 $\sigma' = \sigma 0$ 或 $\sigma' = \sigma 1$，并且 $V(p) = \Setabs{\sigma
    0}{\sigma \in \Bin^*}$ 和 $V(q) = \Setabs{\sigma 1}{\sigma \in
    \Bin^* \setminus \{1\}}$。下图展示了该模型：
  \begin{center}
    \begin{tikzpicture}[modal,
        node distance=2cm
      ]

      \node[world] (0) [label={below:\mTrue{p}\\\mFalse{q}},font=\tiny] {$0$};
      \node[world] (00) [label={below:\mTrue{p}\\\mFalse{q}},
        above right=of 0,font=\tiny] {$00$};
      \node[world] (000) [label={below:\mTrue{p}\\\mFalse{q}},
        right=of 00, yshift=.8cm, font=\tiny] {$000$};
      \node[phantom] (0000) [right=of 000, yshift=.5cm] {};
      \node[phantom] (0001) [right=of 000, yshift=-.5cm] {};
      \node[world] (001) [label={below:\mFalse{p}\\\mTrue{q}},
        right=of 00, yshift=-.8cm, font=\tiny] {$001$};
      \node[phantom](0010) [right=of 001, yshift=.5cm] {};
      \node[phantom](0011) [right=of 001, yshift=-.5cm] {};
      \node[world] (01) [label={below:\mFalse{p}\\\mTrue{q}},
        below right=of 0,font=\tiny] {$01$};
      \node[world] (010) [label={below:\mTrue{p}\\\mFalse{q}},
        right=of 01, yshift=.8cm,font=\tiny] {$010$};
      \node[phantom] (0100) [right=of 010, yshift=.5cm] {};
      \node[phantom] (0101) [right=of 010, yshift=-.5cm] {};
      \node[world] (011) [label={below:\mFalse{p}\\\mTrue{q}},
        right=of 01, yshift=-.8cm,font=\tiny] {$011$};
      \node[phantom] (0110) [right=of 011, yshift=.5cm] {};
      \node[phantom] (0111) [right=of 011, yshift=-.5cm] {};

      \draw[->] (0) to (00);
      \draw[->] (0) to (01);
      \draw[->] (00) to (000);
      \draw[dotted] (000) to (0000);
      \draw[dotted] (000) to (0001);
      \draw[->] (00) to (001);
      \draw[dotted] (001) to (0010);
      \draw[dotted] (001) to (0011);
      \draw[->] (01) to (010);
      \draw[dotted] (010) to (0100);
      \draw[dotted] (010) to (0101);
      \draw[->] (01) to (011);
      \draw[dotted] (011) to (0110);
      \draw[dotted] (011) to (0111);
    \end{tikzpicture}
  \end{center}
  对于每个世界~$w$，我们都有 $\mSat/{M}{\Box(p \lor q) \lif (\Box p \lor \Box q)}[w]$。

  设 $\Gamma$ 为公式~$\Box(p \lor q) \lif
  (\Box p \lor \Box q)$ 的所有子公式的集合。$W^*$ 和~$V^*$ 分别是什么？模型~$\mModel{M}$ 的最细过滤的可达关系是什么？最粗过滤的呢？
\end{prob}

\end{document}